\documentclass[aspectratio=169, 12pt]{beamer}

\usepackage[T2A]{fontenc}
\usepackage[utf8]{inputenc}
\usepackage[russian]{babel}
\usepackage{tikz}
\usetikzlibrary{positioning, arrows.meta, shapes.geometric, shapes.multipart, shapes.symbols, fit, calc}
\usepackage{booktabs}
\usepackage{array}

% --- Тема и цвета ---
\usetheme{Madrid}
\usecolortheme{whale}

\definecolor{accent}{HTML}{4A9EFF}
\definecolor{danger}{HTML}{FF6B6B}
\definecolor{warn}{HTML}{FFA94D}
\definecolor{ok}{HTML}{51CF66}
\definecolor{purple}{HTML}{BE4BDB}
\definecolor{darkbg}{HTML}{2D3436}

\setbeamercolor{frametitle}{bg=darkbg, fg=white}
\setbeamercolor{title}{bg=darkbg, fg=white}
\setbeamercolor{block title}{bg=accent, fg=white}
\setbeamertemplate{navigation symbols}{}

% --- TikZ стили (вне фреймов, чтобы # работал) ---
\tikzset{
    arr/.style={-{Stealth[length=2mm]}, thick, gray!70},
    darr/.style={-{Stealth[length=2mm]}, thick, dashed, gray!50},
    baccent/.style={draw, rounded corners, minimum height=0.8cm, minimum width=2cm, font=\small, fill=accent!15, draw=accent!60, align=center},
    bwarn/.style={draw, rounded corners, minimum height=0.8cm, minimum width=2cm, font=\small, fill=warn!15, draw=warn!60, align=center},
    bok/.style={draw, rounded corners, minimum height=0.8cm, minimum width=2cm, font=\small, fill=ok!15, draw=ok!60, align=center},
    bdanger/.style={draw, rounded corners, minimum height=0.8cm, minimum width=2cm, font=\small, fill=danger!15, draw=danger!60, align=center},
    bpurple/.style={draw, rounded corners, minimum height=0.8cm, minimum width=2cm, font=\small, fill=purple!15, draw=purple!60, align=center},
    saccent/.style={draw, rounded corners=3pt, minimum height=0.65cm, minimum width=1.7cm, font=\scriptsize, fill=accent!15, draw=accent!50, text width=1.5cm, align=center},
    swarn/.style={draw, rounded corners=3pt, minimum height=0.65cm, minimum width=1.7cm, font=\scriptsize, fill=warn!15, draw=warn!50, text width=1.5cm, align=center},
    sok/.style={draw, rounded corners=3pt, minimum height=0.65cm, minimum width=1.7cm, font=\scriptsize, fill=ok!15, draw=ok!50, text width=1.5cm, align=center},
    sdanger/.style={draw, rounded corners=3pt, minimum height=0.65cm, minimum width=1.7cm, font=\scriptsize, fill=danger!15, draw=danger!50, text width=1.5cm, align=center},
    spurple/.style={draw, rounded corners=3pt, minimum height=0.65cm, minimum width=1.7cm, font=\scriptsize, fill=purple!15, draw=purple!50, text width=1.5cm, align=center},
    dbstyle/.style={cylinder, shape border rotate=90, draw=gray!60, fill=gray!10, minimum height=0.9cm, minimum width=1.6cm, font=\small, align=center, aspect=0.3},
    dbpurple/.style={cylinder, shape border rotate=90, draw=purple!60, fill=purple!15, minimum height=0.9cm, minimum width=1.6cm, font=\small, align=center, aspect=0.3},
    cloudstyle/.style={cloud, cloud puffs=10, cloud puff arc=120, draw=danger!50, fill=danger!8, minimum width=2.2cm, minimum height=0.9cm, font=\small, align=center},
    cloudwarn/.style={cloud, cloud puffs=10, cloud puff arc=120, draw=warn!50, fill=warn!8, minimum width=2.8cm, minimum height=0.9cm, font=\scriptsize, align=center},
    docstyle/.style={tape, draw=ok!60, fill=ok!15, minimum height=0.9cm, minimum width=1.8cm, font=\small, align=center, tape bend top=none},
}

\title{n9n Workflow Editor}
\subtitle{AI-конструктор агентов с экспортом в код}
\date{2026}

\begin{document}

% ===== ТИТУЛ =====
\begin{frame}
\titlepage
\end{frame}

% ===== БЛОК 1: n8n — история =====
\begin{frame}{Что такое n8n}
\begin{columns}[T]
\begin{column}{0.5\textwidth}
\textbf{Факты}
\begin{itemize}
    \item 2019 --- Jan Oberhauser, Берлин
    \item n8n = \textbf{nodemation} (node + automation)
    \item Open-source, self-hosted
    \item Визуальная автоматизация без кода
    \item Альтернатива Zapier / Make
\end{itemize}
\end{column}
\begin{column}{0.5\textwidth}
\textbf{Рост после AI-бума}
\begin{itemize}
    \item 147K+ звёзд GitHub
    \item 200K+ пользователей
    \item AI-ноды: Agent, LangChain, OpenAI
    \item Выручка $\times$5 после AI-пивота
    \item Series C --- \$180M, оценка \$2.5B
\end{itemize}
\end{column}
\end{columns}
\end{frame}

% ===== БЛОК 2: Архитектура n8n =====
\begin{frame}{Архитектура n8n --- Single Mode}
\centering
\begin{tikzpicture}[node distance=0.5cm and 0.8cm]
    \node[baccent] (ui) {Editor UI};
    \node[baccent, right=of ui] (api) {REST API};
    \node[baccent, right=of api] (sched) {Scheduler};
    \node[bwarn, below=0.7cm of api] (engine) {Execution Engine};
    \node[bpurple, below=0.7cm of engine] (nodes) {Node System (400+)};

    \node[draw=gray, dashed, rounded corners, fit=(ui)(api)(sched)(engine)(nodes), inner sep=0.3cm, label={[font=\footnotesize, gray]above:Один процесс Node.js}] {};

    \node[dbstyle, right=3.1cm of engine] (db) {SQLite / PG};
    \node[bdanger, right=2.7cm of nodes] (ext) {Внешние API};

    \draw[arr] (ui) -- (api);
    \draw[arr] (api) -- (engine);
    \draw[arr] (sched) |- (engine);
    \draw[arr] (engine) -- (nodes);
    \draw[arr] (engine) -- (db);
    \draw[arr] (nodes) -- (ext);
\end{tikzpicture}
\end{frame}

% ===== Проблемы n8n =====
\begin{frame}{Проблемы n8n под нагрузкой}
\centering
\renewcommand{\arraystretch}{1.4}
\begin{tabular}{>{\bfseries\color{danger}}l l}
\toprule
Проблема & Суть \\
\midrule
Однопоточность & Тяжёлая нода блокирует всё, включая UI \\
Sequential & Ноды --- последовательно, не параллельно \\
Один процесс & UI + scheduler + engine делят ресурсы \\
31\% failures & 10 вебхуков $\rightarrow$ треть запросов падает \\
Масштабирование & Queue Mode + Redis + K8s --- отдельная задача \\
\bottomrule
\end{tabular}
\end{frame}

% ===== Queue Mode =====
\begin{frame}{n8n Queue Mode --- масштабирование}
\centering
\begin{tikzpicture}[node distance=0.6cm and 1.5cm]
    \node[baccent] (main) {Main Instance};
    \node[dbstyle, right=of main, fill=warn!15, draw=warn!60] (redis) {Redis};
    \node[bok, right=of redis, yshift=1.1cm] (w1) {Worker 1};
    \node[bok, right=of redis, yshift=0.0cm] (w2) {Worker 2};
    \node[font=\large, gray, right=of redis, yshift=-0.7cm] (dots) {\dots};
    \node[bok, right=of redis, yshift=-1.4cm] (w3) {Worker N};
    \node[dbpurple, right=of w2] (db) {PostgreSQL};

    \draw[arr] (main) -- (redis);
    \draw[arr] (redis) -- (w1);
    \draw[arr] (redis) -- (w2);
    \draw[arr] (redis) -- (w3);
    \draw[arr] (w1) -- (db);
    \draw[arr] (w2) -- (db);
    \draw[arr] (w3) -- (db);
\end{tikzpicture}

\vspace{0.5cm}
\textcolor{gray}{Решает проблемы, но это уже \textbf{инфраструктурная задача} --- далеко от <<no-code>>}
\end{frame}

% ===== БЛОК 3: Т-Банк =====
\begin{frame}{Т-Банк и <<эмализация>>}
\begin{itemize}
    \item 10\,000+ IT-специалистов, собственный AI-центр
    \item \textbf{<<Эмализация>>} --- внедрение AI во все отделы
    \item Популярная задача --- \textbf{создание нейроагента}
    \item За задачу часто берутся начинающие аналитики
    \item Используют корпоративный self-hosted n8n
\end{itemize}
\end{frame}

% ===== Пайплайн создания агента =====
\begin{frame}{Пайплайн создания агента сейчас}
\centering
\begin{tikzpicture}[node distance=0.15cm]
    \node[saccent] (a1) {Аналитик\\получает задачу};
    \node[saccent, right=of a1] (a2) {Читает\\доки n8n};
    \node[saccent, right=of a2] (a3) {Собирает\\из нод};
    \node[sok, right=of a3] (a4) {Тестирует};
    \node[sdanger, right=of a4] (a5) {Нагрузка\\растёт};
    \node[swarn, right=of a5] (a6) {Разработчик\\переписывает};
    \node[sok, right=of a6] (a7) {Деплой};

    \draw[arr] (a1) -- (a2);
    \draw[arr] (a2) -- (a3);
    \draw[arr] (a3) -- (a4);
    \draw[arr] (a4) -- (a5);
    \draw[arr] (a5) -- (a6);
    \draw[arr] (a6) -- (a7);
\end{tikzpicture}

\vspace{0.5cm}
\begin{columns}[T]
\begin{column}{0.5\textwidth}
\textcolor{danger}{\textbf{Проблемы:}}
\begin{itemize}\small
    \item Аналитик тратит время на доки
    \item Перенос в код --- ручная работа
    \item Нужен второй человек
\end{itemize}
\end{column}
\begin{column}{0.5\textwidth}
\textcolor{danger}{\textbf{Итого:}}
\begin{itemize}\small
    \item \textbf{$\sim$1 месяц}
    \item \textbf{2 человека}
    \item Куча ручной работы
\end{itemize}
\end{column}
\end{columns}
\end{frame}

% ===== БЛОК 3.5: Все агенты одинаковы =====
\begin{frame}{Все агенты устроены одинаково}
\centering
\begin{tikzpicture}[node distance=1.5cm]
    \node[baccent, minimum width=2.2cm] (in) {Вход\\{\scriptsize сообщение / вебхук}};
    \node[docstyle, above right=0.5cm and 1.5cm of in, minimum width=2.2cm] (ctx) {Контекст\\{\scriptsize база знаний / API}};
    \node[bpurple, right=2.5cm of in, minimum width=2.2cm] (llm) {LLM\\{\scriptsize модель + промпт}};
    \node[bwarn, right=2.5cm of llm, minimum width=2.2cm] (out) {Выход\\{\scriptsize ответ / API}};

    \draw[arr] (in) -- (llm);
    \draw[darr] (ctx) -- (llm);
    \draw[arr] (llm) -- (out);
\end{tikzpicture}

\vspace{0.3cm}
\renewcommand{\arraystretch}{1.15}
\small
\begin{tabular}{llll}
\toprule
\textbf{Агент} & \textbf{Контекст} & \textbf{Промпт} & \textbf{Выход} \\
\midrule
Саппорт-бот & FAQ, доки & <<Оператор поддержки>> & Чат \\
HR-ассистент & Регламенты & <<HR-консультант>> & Мессенджер \\
Аналитик & SQL, дашборды & <<Data-аналитик>> & Отчёт \\
Code-reviewer & Стайлгайд & <<Ревьюер кода>> & PR \\
\bottomrule
\end{tabular}

\vspace{0.2cm}
\textcolor{gray}{\small Одна архитектура $\rightarrow$ можно \textbf{стандартизировать}, \textbf{генерировать} и \textbf{конвертировать в код}}
\end{frame}

% ===== БЛОК 4: Архитектура n9n =====
\begin{frame}{n9n --- архитектура}
\centering
\begin{tikzpicture}[node distance=0.35cm and 0.5cm]
    \node[spurple, minimum width=2cm] (land) {Landing\\+ Chat};
    \node[spurple, right=of land, minimum width=2cm] (editor) {Flow Editor};
    \node[spurple, right=of editor, minimum width=2cm] (export_ui) {Export + Bot};

    \node[draw=purple!40, dashed, rounded corners, fit=(land)(editor)(export_ui), inner sep=0.25cm, label={[font=\scriptsize, purple]above:Frontend (React + Vite)}] {};

    \node[saccent, below=1.3cm of land, minimum width=1.7cm] (chat) {chat};
    \node[saccent, right=of chat, minimum width=1.7cm] (gen) {generate};
    \node[saccent, right=of gen, minimum width=1.7cm] (exec) {run};
    \node[saccent, right=of exec, minimum width=1.7cm] (exp) {export};
    \node[saccent, right=of exp, minimum width=1.7cm] (bot) {run-bot};

    \node[draw=accent!40, dashed, rounded corners, fit=(chat)(gen)(exec)(exp)(bot), inner sep=0.25cm, label={[font=\scriptsize, accent]below:Backend (FastAPI)}] {};

    \node[swarn, below=1.3cm of exec, minimum width=3cm] (llm) {DeepSeek / OpenAI / Ollama};

    \draw[arr] (land) -- (chat);
    \draw[arr] (land) -- (gen);
    \draw[arr] (editor) -- (exec);
    \draw[arr] (editor) -- (exp);
    \draw[arr] (export_ui) -- (bot);
    \draw[darr] (gen) -- (llm);
    \draw[darr] (exec) -- (llm);
    \draw[darr] (chat) -- (llm);
\end{tikzpicture}
\end{frame}

% ===== Пайплайн n9n =====
\begin{frame}{Пайплайн создания агента в n9n}
\centering
\begin{tikzpicture}[node distance=0.2cm]
    \node[saccent, minimum width=1.8cm] (s1) {Описал задачу\\в чате};
    \node[saccent, right=of s1, minimum width=1.8cm] (s2) {AI генерирует\\workflow};
    \node[saccent, right=of s2, minimum width=1.8cm] (s3) {Правки\\чат / визуально};
    \node[sok, right=of s3, minimum width=1.8cm] (s4) {Тест в\\интерфейсе};
    \node[sok, right=of s4, minimum width=1.8cm] (s5) {Export\\Python + Docker};
    \node[sok, right=of s5, minimum width=1.8cm] (s6) {Деплой};

    \draw[arr] (s1) -- (s2);
    \draw[arr] (s2) -- (s3);
    \draw[arr] (s3) -- (s4);
    \draw[arr] (s4) -- (s5);
    \draw[arr] (s5) -- (s6);
\end{tikzpicture}

\vspace{0.5cm}
\begin{columns}[T]
\begin{column}{0.5\textwidth}
\textcolor{ok}{\textbf{Результат:}}
\begin{itemize}\small
    \item \textbf{Минуты--часы} вместо месяца
    \item \textbf{1 человек} вместо двух
    \item Код генерируется автоматически
\end{itemize}
\end{column}
\begin{column}{0.5\textwidth}
\textcolor{ok}{\textbf{На выходе:}}
\begin{itemize}\small
    \item \texttt{main.py} --- Telegram-бот
    \item \texttt{workflow.py} --- логика агента
    \item \texttt{Dockerfile} --- контейнер
    \item \texttt{requirements.txt}
\end{itemize}
\end{column}
\end{columns}
\end{frame}

% ===== Типы нод =====
\begin{frame}{Типы нод}
\centering
\begin{tikzpicture}[node distance=1cm and 1.8cm]
    \node[baccent, minimum width=1.8cm] (trig) {Trigger};
    \node[bwarn, right=of trig, yshift=0.6cm, minimum width=1.8cm] (http) {HTTP};
    \node[bdanger, right=of trig, yshift=-0.6cm, minimum width=1.8cm] (cond) {Condition};
    \node[bpurple, right=of http, yshift=-0.6cm, minimum width=1.8cm] (agent) {Agent};
    \node[docstyle, above=0.6cm of agent, minimum width=1.8cm] (know) {Knowledge};

    \draw[arr] (trig) -- (http);
    \draw[arr] (trig) -- (cond);
    \draw[arr] (http) -- (agent);
    \draw[arr] (cond) -- (agent);
    \draw[darr] (know) -- node[right, font=\scriptsize, gray]{context} (agent);
\end{tikzpicture}

\vspace{0.4cm}
\small
\begin{tabular}{ll}
\toprule
\textbf{Нода} & \textbf{Назначение} \\
\midrule
\textcolor{accent}{Trigger} & Точка входа (manual, webhook) \\
\textcolor{warn}{HTTP} & Вызов внешнего API \\
\textcolor{danger}{Condition} & Проверка (empty, not\_empty, contains) \\
\textcolor{ok}{Knowledge} & RAG-контекст (документы, URL) \\
\textcolor{purple}{Agent} & LLM-вызов (DeepSeek, OpenAI, Ollama) \\
\bottomrule
\end{tabular}
\end{frame}

% ===== Сравнение =====
\begin{frame}{n8n vs n9n}
\centering
\renewcommand{\arraystretch}{1.4}
\small
\begin{tabular}{l c c}
\toprule
& \textbf{n8n} & \textbf{n9n (наш)} \\
\midrule
Создание & Руками + доки & AI генерирует \\
Редактирование & Drag \& drop & Чат + визуально \\
Масштабирование & Queue Mode, K8s & Export $\rightarrow$ Python + Docker \\
Перенос в код & Ручной, разработчик & Автоматический \\
Людей & 2 (аналитик + dev) & 1 (аналитик) \\
Время & $\sim$1 месяц & Минуты--часы \\
Запуск бота & Внешний деплой & Одна кнопка \\
\bottomrule
\end{tabular}
\end{frame}

% ===== БЛОК 5: Будущее =====
\begin{frame}{Что дальше}
\begin{columns}[T]
\begin{column}{0.5\textwidth}
\textbf{\textcolor{accent}{Ближайшее}}
\begin{itemize}\small
    \item Новые ноды: database, transform, loop
    \item Визуальный дебаггер
    \item Шаблоны workflow
    \item История версий
\end{itemize}

\vspace{0.3cm}
\textbf{\textcolor{purple}{Продвинутое}}
\begin{itemize}\small
    \item Multi-agent orchestration
    \item RAG-пайплайн (chunking + vector store)
    \item Memory / контекст диалога
    \item Мониторинг (latency, tokens, ошибки)
\end{itemize}
\end{column}
\begin{column}{0.5\textwidth}
\textbf{\textcolor{ok}{Инфраструктура}}
\begin{itemize}\small
    \item One-click deploy (Railway, VPS)
    \item CI/CD-интеграция
    \item Persistent storage (PostgreSQL)
    \item Auth, multi-user
\end{itemize}

\vspace{0.3cm}
\textbf{\textcolor{warn}{Интеграции}}
\begin{itemize}\small
    \item MCP (Model Context Protocol)
    \item Jira, Confluence, Slack, Bitrix24
    \item GigaChat, YandexGPT, T-lite
\end{itemize}
\end{column}
\end{columns}
\end{frame}

% ===== ИТОГ =====
\begin{frame}{Итог}
\centering
\vspace{0.4cm}
{\normalsize Мы не решаем одну конкретную боль и не автоматизируем одну рутину.}

\vspace{0.5cm}
{\large\bfseries Мы даём \textcolor{accent}{полноценный инструмент},\\[0.2cm]
с помощью которого можно автоматизировать\\[0.15cm]
\textcolor{accent}{тысячи задач} --- быстро и без боли\\[0.15cm]
встроить в IT-ландшафт предприятия любого размера.}
\end{frame}

\end{document}
